Large language models are typically evaluated as answer generators, but real assistants must first decide whether answering is appropriate. We study next-action selection under defective inputs as an evidence-conditioned first-move decision: given a user query and optional evidence, a model must choose whether to answer, ask a clarification, challenge a false or stale premise, or abstain when support is insufficient. We construct a benchmark in which answerable controls, false premises, stale premises, and conflicting evidence share one action-decision record, action ontology, utility-centered metric suite, and completed human-validation queue for the active evidence bundle. On the completed Day-1 local model-and-prompt matrix, open instruction and reasoning checkpoints remain far from saturated. In particular, \DayOneScaleReasoningLargestMatchedReasoningModel{} reaches only \DayOneScaleReasoningLargestMatchedReasoningActionAcc{} action accuracy and \DayOneScaleReasoningLargestMatchedReasoningUtility{} utility on the development split, below the strongest instruction baselines, with weak false-premise interruption and low JSON adherence. A lightweight decision-first prompt improves Qwen2.5-1.5B on the quick+stale split, but negative ablations show that simply adding more caution or critique language can damage answerable behavior. These results suggest that next-action calibration is a distinct capability: fluent answer generation and chain-of-thought reasoning do not guarantee the right first move.
